% ======================================================================
\section{Worked examples: one row of each Group}\label{sec:examples}
% ======================================================================
% Sources: Example 1: app_wavefunctions (calc_main_checked/Calculation_main_1_checked.tex), group_I_checked, limits_checked/Divergent_terms_checked.tex.
% Example 2: group_II_checked/Group_II_Row_IV_checked.tex (all steps). Example 3: group_III_checked/Group_III_checked.tex (Row V),
% evolution_checked/BFKL_checked.tex, Other_terms_checked.tex (split of 1/p+), final_explicit/Final_terms_explicit_2_fixed.tex.
% The intermediate algebra quoted here is re-checked in scripts/check_examples.py.

\purpose{The fifteen rows are evaluated by the same sequence of operations, and the remaining sections quote their results. To make these results reproducible we follow here three rows in full, from the coefficients of $\Omega$ to the form in which they appear in the final result: the one-$\rho$ part of Row~I.II (a virtual correction with an ultraviolet pole), Row~II.IV (the real emission with the collinear divergence, which produces the DGLAP term) and Row~III.V (an interference term cubic in the charges, with a recoil phase, a plus prescription and an instantaneous term). Between them the three examples contain every type of longitudinal and transverse integral, every kind of subtraction, and the treatment of the ordered charges that occur in the paper. The other rows are stated in Sections~\ref{sec:rapidity}--\ref{sec:finite} and in Appendix~\ref{app:finite}.}

\subsection{The common procedure}\label{sec:ex_procedure}

Every row is evaluated in six steps.
\begin{enumerate}
\item \emph{Insert the coefficients} of Section~\ref{sec:coefficients} into the product of amplitudes that defines the row (Table~\ref{tab:rows}), with the conjugated factor to the left and at the primed points.
\item \emph{Colour algebra.} Wilson lines at the same point are combined with $U^TU=1$ and $f^{abc}U^{bd}U^{ce}=U^{ah}f^{hde}$. The charges are not moved.
\item \emph{Transverse $\delta$-functions and polarisation sums.} The $\delta$-function of $\mf C_1$ puts the parent gluon at the $p^+$-weighted centre of its daughters; this is the origin of all recoil phases. Sums over the polarisations of unobserved gluons are done, in $d_\perp$ dimensions where they multiply a pole.
\item \emph{Longitudinal integral.} The $p^+$ integral is written in terms of a momentum fraction $\xi$. Its end points give the rapidity logarithms; these are taken out by plus prescriptions, and the subtracted terms are recorded in the BFKL and JIMWLK lists.
\item \emph{Transverse integrals} that can be done in closed form, in $d=2-\epsilon$ dimensions if they diverge.
\item \emph{Limits and final form.} $\Lambda\to0$, $\vee\to\infty$ in the subtracted remainder; for a cubic row, decomposition of the ordered product of charges into its symmetrised part and contact terms; average over the projectile.
\end{enumerate}

\subsection{Example 1 (Group I): the one-\texorpdfstring{$\rho$}{rho} part of Row I.II}\label{sec:ex_I_II}

\paragraph{The row.} Row~I.II is the interference of the LO amplitude with the $O(g^3)$ one-gluon component of the wave function,
\begin{align}\label{eq:ex1_row}
\frac{d^3N^{\rm I.II}}{d^3k}=\frac{1}{(2\pi)^3}\Big\langle\mathcal A^{a\dagger}_{{\rm LO}\,i}(k)\int_we^{-i\v k\cdot\v w}\Big[U^{ab}(\v w)\,\mf A^{(3)b}_i(k^+,-\v w)-\bar{\mf A}^{(3)a}_i(k^+,-\v w)\Big]\Big\rangle+{\rm c.c.}
\end{align}
We consider the part of $\mf A^{(3)}$ that is linear in $\rho$. It is the simplest part of the calculation and at the same time the place where the numbers $\frac{11}{3}$ and $\frac{67}{9}-\frac{\pi^2}{3}$ arise.

\paragraph{Step 1: the self-energy and its longitudinal integrals.} The diagram of Figure~\ref{fig:vac_g3_onerho} is a gluon of momentum $(k^+,\v k)$ that fluctuates into a pair with momentum fractions $\xi$ and $1-\xi$ and recombines. After the transverse loop integral in $d=2-\epsilon$ dimensions the one-gluon component of the state is the LO component multiplied by (Appendix~\ref{app:wavefunctions})
\begin{align}\label{eq:ex1_SE}
+\frac{g^2N_c}{8\pi^2}\int_{\delta}^{1-\delta}d\xi\,\Big[-2+\xi(1-\xi)+\frac{1}{\xi(1-\xi)}\Big]\Big[-\frac2\epsilon+\log\frac{\v k^2\,\xi(1-\xi)}{\mu^2}\Big]\,,\qquad\delta=\frac{\Lambda}{k^+}\,.
\end{align}
The first bracket is the splitting kernel: it is identical to $\xi(1-\xi)+\frac{\xi}{1-\xi}+\frac{1-\xi}{\xi}=\mathcal P_{gg+}(\xi)/(2N_c)$, Eq.~\eqref{eq:Pggplus}; the second is the transverse loop, whose ultraviolet pole is independent of $\xi$ and whose finite part depends on $\xi$ through the virtuality $\v k^2\xi(1-\xi)$ of the pair. Both members of the pair are soft modes, hence the limits $\delta$ and $1-\delta$. Three integrals are needed,
\begin{align}\label{eq:ex1_ints}
\int_\delta^{1-\delta}d\xi\Big[-2+\xi(1-\xi)+\frac{1}{\xi(1-\xi)}\Big]&=-\frac{11}{6}-2\log\delta\,,\notag\\
\int_0^1d\xi\,\big[-2+\xi(1-\xi)\big]\log\big(\xi(1-\xi)\big)&=\frac{67}{18}\,,\notag\\
\int_\delta^{1-\delta}d\xi\,\frac{\log(\xi(1-\xi))}{\xi(1-\xi)}&=-\log^2\delta-\frac{\pi^2}{3}\,,
\end{align}
up to terms that vanish with $\delta$. The first gives the coefficient of the pole: $-\frac{11}{6}$ is the integral of the non-singular part of the kernel and is half of $\beta_0/N_c=\frac{11}{3}$, and $-2\log\delta$ comes from the two soft end points. The second and third give the constant. Hence Eq.~\eqref{eq:ex1_SE} is
\begin{align}\label{eq:ex1_SEres}
-\frac{g^2N_c}{16\pi^2}\Big\{\Big(\frac{11}{3}-4\log\frac{k^+}{\Lambda}\Big)\Big[-\frac2\epsilon+\log\frac{\v k^2}{\mu^2}\Big]+2\log^2\frac{k^+}{\Lambda}-\frac{67}{9}+\frac{2\pi^2}{3}\Big\}\,.
\end{align}
\emph{Reading:} $\frac{11}{3}$ multiplies the ultraviolet pole and is the pure-glue $\beta_0/N_c$; $\frac{67}{9}$ is the finite part of the same loop; the double logarithm and $\frac{2\pi^2}{3}$ come from the soft end points of the pair.

\paragraph{Step 2: the other one-$\rho$ contributions.} The commutator parts of the diagrams with two and three valence vertices are also linear in $\rho$. In them the measured gluon carries a fraction $\xi$ of the momentum of a parent that has transferred the rest to the valence charges, and the parent must be a soft mode, $k^+/\xi<\vee$. Their $\xi$ integrals (Appendix~\ref{app:wavefunctions}) change the coefficient of the pole from $\frac{11}{3}-4\log\frac{k^+}{\Lambda}$ to $\frac{11}{3}-2\log\frac{k^+}{\Lambda}+\log(1-\frac{k^+}{\vee})$, cancel the double logarithm of $\Lambda$, and change $\frac{2\pi^2}{3}$ into $\frac{\pi^2}{3}$. The result is Eq.~\eqref{eq:A3onerho}; in the eikonal limit $\vee\to\infty$,
\begin{align}\label{eq:ex1_A3}
\Big[\mf A^{(3)a}_{2i}(k^+,\v k)\Big]_\rho=-\frac{g^2N_c}{16\pi^2}\,\mf A^{(1)a}_{2i}(k^+,\v k)\,\Big\{\Big[-\frac2\epsilon+\log\frac{\v k^2}{\mu^2}\Big]\Big(\frac{11}{3}-2\log\frac{k^+}{\Lambda}\Big)+\frac{\pi^2}{3}-\frac{67}{9}\Big\}\,,
\end{align}
with $\mf A^{(1)a}_{2i}(k^+,\v k)=-\sqrt2\,g\,k^i\rho^a(-\v k)/(\sqrt{k^+}\,\v k^2)$. \emph{Nuance:} the cancellation of $\log^2(k^+/\Lambda)$ between the self-energy and the commutator terms means that the split of the one-$\rho$ part into ``self-energy'' and ``vertex corrections'' is not meaningful term by term: each has a double logarithm of $\Lambda$, and only the sum is free of it.

\paragraph{Step 3: transverse Fourier transform.} The coefficient is needed at the transverse position $\v w$ of the gluon. With Eqs.~\eqref{eq:FT_K} and~\eqref{eq:FT_log},
\begin{align}
\frac{k^i}{\v k^2}\;\to\;\Kv^i(\v x-\v w)\,,\qquad\frac{k^i}{\v k^2}\log\frac{\v k^2}{\mu^2}\;\to\;-\mathcal B(\v x-\v w)\,\Kv^i(\v x-\v w)\,,\qquad\mathcal B(\v r)=\log\frac{\v r^2\mu^2}{c_0^2}\,,
\end{align}
so that the momentum-space logarithm becomes minus the logarithm of the distance between the emitting charge and the gluon, with the constant $c_0=2e^{-\gamma_E}$. This gives Eq.~\eqref{eq:A3onerhoX}: point by point in $\v x$, the one-$\rho$ part of $\mf A^{(3)}$ is the integrand of $\mf A^{(1)}$, Eq.~\eqref{eq:A1}, times $-\frac{g^2N_c}{16\pi^2}\{[-\frac2\epsilon-\mathcal B(\v x-\v w)](\frac{11}{3}-2L)+\frac{\pi^2}{3}-\frac{67}{9}\}$, $L=\log(k^+/\Lambda)$.

\paragraph{Step 4: insertion into the row.} Because of this proportionality the bracket of Eq.~\eqref{eq:ex1_row} has the colour structure of the LO amplitude, $[U(\v w)-U(\v x)]^{ab}\rho^b(\v x)$, and the row is the LO integrand with an additional factor. Including the complex conjugate, and with the notation of Eq.~\eqref{eq:dg_not},
\begin{align}\label{eq:ex1_res}
\frac{d^3N^{\rm I.II}}{d^3k}\bigg|_{1\rho}=\ &\frac{\alpha_sN_c}{4\pi}\,\mathcal N_0\int_{x,w,x',w'}e^{-i\v k\cdot(\v w-\v w')}\,\mathcal O(w,w')\notag\\
&\times\Big\{\Big(\frac{11}{3}-2L\Big)\Big[\frac4\epsilon+\mathcal B(\v x-\v w)+\mathcal B(\v x'-\v w')\Big]+2\Big(\frac{67}{9}-\frac{\pi^2}{3}\Big)\Big\}\,.
\end{align}

\paragraph{Step 5: where the three pieces go.}
\begin{itemize}
\item The term with $\frac{11}{3}$, pole and logarithms, is Eq.~\eqref{eq:dg_II}. It is cancelled exactly by the part of Row~I.V in which the recombining pair is at one point, Eq.~\eqref{eq:dg_V} (Section~\ref{sec:dg_poles}). It does not appear in the final result.
\item The terms with $L=\log(k^+/\Lambda)$ are assigned to the JIMWLK sector, in which the poles proportional to $L$ cancel (Table~\ref{tab:poles}).
\item The constant is the genuine two-$\rho$ finite term of the row, Eq.~\eqref{eq:S-I.II-2o}: with $\frac{\alpha_sN_c}{4\pi}\mathcal N_0=\frac{1}{(2\pi)^3}\frac{g^4N_c}{32\pi^4}\frac{1}{k^+}$ and $\mathcal O(w,w')$ written as $-\Kv(x'-w')\cdot\Kv(x-w)[U(x')-U(w')][U(w)-U(x)]\rho\rho$, the printed coefficient $(-\frac{67}{9}+\frac{\pi^2}{3})$ with ``$+{\rm c.c.}$'' is Eq.~\eqref{eq:ex1_res}.
\end{itemize}
The two- and three-$\rho$ parts of the row are obtained in the same way from Eqs.~\eqref{eq:A3tworho} and~\eqref{eq:A3threerho}; the first gives the cubic part of the row with the momentum-space kernel $\mathfrak K_{\rm I.II}$, the second the four-$\rho$ terms of Section~\ref{sec:fourrho}.

\subsection{Example 2 (Group II): Row II.IV and the origin of the DGLAP term}\label{sec:ex_II_IV}

\paragraph{The row.} Row~II.IV is the square of $\mathcal A_{CA}$, Eq.~\eqref{eq:A2gamp}: a gluon of momentum $p^++k^+$, emitted at LO, splits after the collision into the measured gluon $(k^+,\v w)$ and a gluon $(p^+,\v z)$ that is integrated over. The parent is a soft mode, so that $\Lambda<p^+<\vee-k^+$.

\paragraph{Step 1: coefficients.} With $\mf C_1$ of Eq.~\eqref{eq:C1} and $\mf A$ of Eq.~\eqref{eq:A1} for the parent at $\v y$ (amplitude) and $\v y'$ (conjugate amplitude),
\begin{align}\label{eq:ex2_start}
\frac{d^3N^{\rm II.IV}}{d^3k}=\ &\frac{1}{(2\pi)^3}\,\frac{g^4}{4\pi^4}\int_\Lambda^{\vee-k^+}\frac{dp^+}{2\pi}\,\frac{p^+k^+}{(p^++k^+)^4}\int_{z,w,w',x,x',y,y'}e^{-i\v k\cdot(\v w-\v w')}\notag\\
&\times\Kv^m(w-z)\,\Kv^k(x-y)\,\Kv^{m'}(w'-z)\,\Kv^{k'}(x'-y')\,\,\tau^{km}_{ij}\,\tau^{k'm'}_{ij}\notag\\
&\times\delta^{(2)}\Big[\v y-\v w+\frac{p^+(\v w-\v z)}{p^++k^+}\Big]\,\delta^{(2)}\Big[\v y'-\v w'+\frac{p^+(\v w'-\v z)}{p^++k^+}\Big]\notag\\
&\times f^{c'd'a}\big[U^{c'e'}(y')-U^{c'e'}(x')\big]U^{d'b}(z)\notag\\
&\times f^{cda}\big[U^{ce}(y)-U^{ce}(x)\big]U^{db}(z)\;\rho^{e'}(x')\rho^e(x)\,,\notag\\
\tau^{km}_{ij}\equiv\ &\delta_{km}\delta_{ij}-\frac{p^++k^+}{p^+}\,\delta_{jm}\delta_{ik}-\frac{p^++k^+}{k^+}\,\delta_{im}\delta_{kj}\,,
\end{align}
where $i$ and $j$ are the polarisations of the two final-state gluons and $k$, $k'$ those of the parent.

\paragraph{Step 2: colour.} The unobserved gluon carries the same Wilson line $U(\v z)$ in the amplitude and in the conjugate amplitude: $U^{d'b}(z)U^{db}(z)=\delta^{d'd}$, and then $f^{c'da}f^{cda}=N_c\delta^{c'c}$. The colour factor becomes $N_c\,[U(y')-U(x')]^{ce'}[U(y)-U(x)]^{ce}\rho^{e'}(x')\rho^e(x)$: $N_c$ times the LO colour structure of the parent. \emph{This is the statement that a final-state splitting does not change the interaction with the target.}

\paragraph{Step 3: momentum fraction, $\delta$-functions, polarisation sum.} With $\xi=k^+/(p^++k^+)$, Eq.~\eqref{eq:xi_defIV}, one has $dp^+\,p^+k^+/(p^++k^+)^4=d\xi\,\xi(1-\xi)/k^+$, the range becomes $k^+/\vee<\xi<k^+/(k^++\Lambda)$, and the $\delta$-functions put the parent at
\begin{align}\label{eq:ex2_parent}
\v y=\v z+\xi(\v w-\v z)\,,\qquad\v y'=\v z+\xi(\v w'-\v z)\,,
\end{align}
the $p^+$-weighted centre of the two daughters. The sum over the polarisations $i,j$ is, in $d_\perp$ dimensions,
\begin{align}\label{eq:ex2_pol}
\xi(1-\xi)\,\tau^{km}_{ij}\,\tau^{k'm'}_{ij}=\ &d_\perp\,\xi(1-\xi)\,\delta_{km}\delta_{k'm'}+2\big(\delta_{mk'}\delta_{km'}-\delta_{km}\delta_{k'm'}\big)\notag\\&+\Big(\frac{\xi}{1-\xi}+\frac{1-\xi}{\xi}\Big)\delta_{mm'}\delta_{kk'}\,.
\end{align}

\paragraph{Step 4: the collinear integral.} In the variables $\v r=\v w-\v z$ (separation of the daughters), $\v u=\v w'-\v w$ and $\v Y=\frac12(\v y+\v y')$, with $\v y=\v Y-\frac\xi2\v u$, $\v y'=\v Y+\frac\xi2\v u$ and unit Jacobian, the only dependence on $\v r$ is in the two splitting kernels. Their integral diverges logarithmically at large $|\v r|$, where the daughters are far apart in the transverse plane, i.e.\ have small relative transverse momentum: this is the collinear divergence. In $d=2-\epsilon$ dimensions
\begin{align}\label{eq:ex2_rint}
\int d^dr\,\frac{r^m\,(r+u)^{m'}}{\v r^2\,(\v r+\v u)^2}
=-\frac\pi2\Big[\frac2\epsilon+2\gamma_E+\log\frac{\v u^2\mu^2}{4}\Big]\delta^{mm'}-\pi\,\frac{u^mu^{m'}}{\v u^2}+O(\epsilon)\,.
\end{align}
(The integral is the Fourier transform of $q^mq^{m'}/\v q^4$; it diverges at $\v q\to0$ and exists for $\epsilon<0$.) \emph{Nuance:} the contraction of Eq.~\eqref{eq:ex2_rint} with Eq.~\eqref{eq:ex2_pol} must be done with $d_\perp=\delta_{mm}=2-\epsilon$, because $d_\perp$ multiplies the pole. The product $-\epsilon\times\frac1\epsilon$ combines with the term $u^mu^{m'}/\v u^2$ into a traceless tensor; with $d_\perp=2$ the bracket $[\frac2\epsilon+\dots]$ would instead be $[\frac2\epsilon+\dots+1]$. The result is
\begin{align}\label{eq:ex2_contr}
\frac{d^3N^{\rm II.IV}}{d^3k}=\ &-\frac{1}{(2\pi)^3}\,\frac{g^4N_c}{8\pi^4}\,\frac{1}{k^+}\int_{k^+/\vee}^{k^+/(k^++\Lambda)}d\xi\int_{u,Y,x,x'}e^{i\v k\cdot\v u}\notag\\
&\times\Big\{\Big[\xi(1-\xi)+\frac{\xi}{1-\xi}+\frac{1-\xi}{\xi}\Big]\Big[\frac2\epsilon+2\gamma_E+\log\frac{\v u^2\mu^2}{4}\Big]\Kv(x-y)\cdot\Kv(x'-y')\notag\\
&\quad+2\,\xi(1-\xi)\Big[\frac{\v u\cdot\Kv(x-y)\;\v u\cdot\Kv(x'-y')}{\v u^2}-\frac12\,\Kv(x-y)\cdot\Kv(x'-y')\Big]\Big\}\notag\\
&\times\big[U(y')-U(x')\big]^{ce'}\big[U(y)-U(x)\big]^{ce}\rho^{e'}(x')\rho^e(x)\,.
\end{align}
The first bracket is the real gluon splitting function $\mathcal P_{gg+}/(2N_c)$, Eq.~\eqref{eq:Pggplus}, before any prescription. It has emerged from the polarisation sum of the light-cone vertices, not been put in.

\paragraph{Step 5: the two end points.} The $\xi$ integral diverges logarithmically at both ends as the cut-offs are removed, and the two divergences have different physical meaning.

\emph{$\xi\to1$: the integrated gluon is soft.} Then $\v y\to\v w$, $\v y'\to\v w'$: the parent is the measured gluon. With $\int d\xi/(1-\xi)=\log\frac{k^+}{\Lambda}+\log(1-\frac{k^+}{\vee})+O(\Lambda/k^+)$ this end point gives $[L+\ell_\vee]$ times the LO structure with the collinear bracket: a contribution to the JIMWLK list, the emission of a soft gluon by the measured gluon after the collision.

\emph{$\xi\to0$: the measured gluon is soft compared with its parent.} Then $\v y,\v y'\to\v Y$: the parent no longer distinguishes the amplitude from the conjugate amplitude, and the colour structure becomes independent of $\v u$. The $\v u$ integral can be done,
\begin{align}\label{eq:ex2_uint}
\int d^2u\,e^{i\v k\cdot\v u}\Big[\frac2\epsilon+2\gamma_E+\log\frac{\v u^2\mu^2}{4}\Big]=-\frac{4\pi}{\v k^2}\qquad(\v k\neq0)\,,
\end{align}
the constant terms, including the pole, giving $\delta^{(2)}(\v k)$. With $\int d\xi/\xi=\log\frac{\vee}{k^+}+O(\Lambda/k^+)$ this end point gives $+2N_c\,\mathcal C\,\frac{4\pi}{\v k^2}\log\frac{\vee}{k^+}\,\mathcal X$ with $\mathcal X$ of Eq.~\eqref{eq:dg_X}: a contribution to the BFKL list. It has no pole. Physically the parent is a fast gluon of the projectile at $\v Y=\v z$ which radiates the measured gluon; the measured gluon sees it as an additional colour charge.

Subtracting and adding the two end-point terms, and changing in the remainder the variable $\v u$ to the separation of the parents, $\v y'-\v y=\xi\v u$, which gives a factor $\xi^{-2}$, the phase $e^{-i\v k\cdot(\v y-\v y')/\xi}$ and the argument $(\v y'-\v y)^2/\xi^2$ of the logarithm, the row becomes
\begin{align}\label{eq:ex2_split}
\frac{d^3N^{\rm II.IV}}{d^3k}=\ &-\mathfrak a\,\big[L+\ell_\vee\big]\,\mathcal N_0\int_{x,y,x',y'}e^{-i\v k\cdot(\v y-\v y')}\Big[\frac2\epsilon+\mathcal B(\v y'-\v y)\Big]\mathcal O(y,y')\notag\\
&+2N_c\,\mathcal C\,\frac{4\pi}{\v k^2}\,\log\frac{\vee}{k^+}\,\mathcal X\notag\\
&-\mathcal C\int_{k^+/\vee}^1d\xi\,\Big\{\frac{\mathcal P_{gg+}(\xi)}{\xi^2}\int_{x,y,x',y'}e^{-i\v k\cdot(\v y-\v y')/\xi}\Big[\frac2\epsilon+\log\frac{(\v y'-\v y)^2\mu^2}{c_0^2\,\xi^2}\Big]\mathcal O(y,y')\notag\\
&\hspace{2.6cm}+\frac{2N_c}{\xi}\,\frac{4\pi}{\v k^2}\,\mathcal X\Big\}\notag\\
&+\frac{d^3N^{\rm II.IV}_{\rm tensor}}{d^3k}\,,
\end{align}
where now $\mathcal P_{gg+}$ carries the plus prescription at $\xi\to1$, Eq.~\eqref{eq:plusxi}, $\mathfrak a=\alpha_sN_c/\pi$, and the last term is the tensor term of Eq.~\eqref{eq:ex2_contr}, which has no pole and is the first term of Eq.~\eqref{eq:S-II.IVrem}. The upper limit has been set to one, at the cost of $O(\Lambda/k^+)$; the lower limit $k^+/\vee$ is kept, because the integrand is not small there before the integral over $\v y-\v y'$ is done.

\paragraph{Step 6: from Eq.~\eqref{eq:ex2_split} to the DGLAP, running-coupling and finite terms.} The second line of Eq.~\eqref{eq:ex2_split} is the collinear term. Three things remain to be done with it, each of which is discussed in Section~\ref{sec:dglap_rc}: its pole is absorbed into the fragmentation function, with a counterterm whose Jacobian is $\xi^{-d}$ and not $\xi^{-2}$ (Section~\ref{sec:dg_msbar}); the term $\frac{\beta_0}{2}\delta(1-\xi)$, which a real emission cannot produce, is added to $\mathcal P_{gg+}$ and subtracted again, the subtracted term being the running-coupling term~\eqref{eq:RC}; and the subtraction at $\xi\to0$ is written as a plus prescription inside the splitting function, Eq.~\eqref{eq:DGLAP}. The last step leaves a finite term: since
\begin{align}\label{eq:ex2_rem}
\int_{k^+/\vee}^1d\xi\,\frac{1-\xi}{\xi}=\log\frac{\vee}{k^+}-1+\frac{k^+}{\vee}\,,
\end{align}
and not $\log(\vee/k^+)$, the plus prescription on $\frac{1-\xi}{\xi}$ subtracts the BFKL term of the row \emph{and} a term with the coefficient $-1$; the latter is the remnant~\eqref{eq:remnant}, $-2N_c\,\mathcal C\frac{4\pi}{\v k^2}\mathcal X$, the second term of Eq.~\eqref{eq:S-II.IVrem}.

\paragraph{What the example shows.} One row has produced all five types of terms of the result: a JIMWLK term ($\xi\to1$), a BFKL term ($\xi\to0$), the DGLAP term (collinear pole between the end points), the running-coupling term (through the $\delta(1-\xi)$ that completes the splitting function) and finite terms (tensor term and remnant). The assignment of each piece is dictated by the region of $\xi$ and of the transverse integral from which it comes.

\subsection{Example 3 (Group III): Row III.V}\label{sec:ex_III_V}

\paragraph{The row.} Row~III.V is $-\mathcal A^\dagger_{CA}\mathcal A^{(1)}_{AA}+{\rm c.c.}$: in the conjugate amplitude the two final-state gluons are the daughters of a LO gluon that has split after the collision; in the amplitude they are two independently emitted gluons, one before and one after the collision. It is cubic in the charges: one charge at $\v x'$ from the conjugate amplitude, two charges at $\v x$ and $\v y$ from the amplitude. The parent gluon of the splitting is a soft mode, $\Lambda<p^+<\vee-k^+$.

\paragraph{Step 1: coefficients.} Inserting $\mf C_1$ and $\mf A$,
\begin{align}\label{eq:ex3_start}
\frac{d^3N^{\rm III.V}}{d^3k}=\ &\frac{1}{(2\pi)^3}\,\frac{ig^4f^{c'd'a}}{4\pi^4}\,\frac{1}{k^{+2}}\int_\Lambda^{\vee-k^+}\frac{dp^+}{2\pi}\int_{w,w',y',z,y,x,x'}e^{-i\v k\cdot(\v w-\v w')}\notag\\
&\times\Kv^m(w'-z)\,\Kv^{k'}(x'-y')\,\Kv^i(y-w)\,\Kv^j(x-z)\;\,\delta^{(2)}\Big[\v w'-\frac{(p^++k^+)\v y'-p^+\v z}{k^+}\Big]\notag\\
&\times\Big[\delta_{k'm}\delta_{ij}-\frac{p^++k^+}{p^+}\,\delta_{jm}\delta_{ik'}-\frac{p^++k^+}{k^+}\,\delta_{im}\delta_{jk'}\Big]\notag\\
&\times\big[U^{c'e'}(y')-U^{c'e'}(x')\big]U^{d'b}(z)\,\rho^{e'}(x')\notag\\
&\times\big[U^{ad}(y)-U^{ad}(w)\big]U^{ce}(x)\,U^{cb}(z)\,\rho^e(x)\rho^d(y)+{\rm c.c.}
\end{align}
The single structure constant is that of $\mf C_1$; the amplitude contains none, since its two emissions are independent.

\paragraph{Step 2: colour.} Again $U^{d'b}(z)U^{cb}(z)=\delta^{d'c}$: the integrated gluon does not see the target in the product. The colour operator is
\begin{align}\label{eq:ex3_colour}
f^{mkl}\big[U^{ma}(y')-U^{ma}(x')\big]\,U^{kb}(x)\,\big[U^{lc}(y)-U^{lc}(w)\big]\;\rho^a(x')\,\rho^b(x)\,\rho^c(y)\,,
\end{align}
of the form $[U-U]\times U\times[U-U]\times f$ that is common to all cubic rows (Section~\ref{sec:fin_singlet}). The charges stand in the order conjugate amplitude, then amplitude, and within the amplitude in the order of $\mathcal A^{(1)}_{AA}$; they are not rearranged.

\paragraph{Step 3: recoil.} The $\delta$-function fixes $\v w'=\v y'+\frac{p^+}{k^+}(\v y'-\v z)$: the measured daughter is displaced from the parent, away from the other daughter. The integral over $\v w'$ therefore replaces the phase by
\begin{align}\label{eq:ex3_phase}
e^{-i\v k\cdot(\v w-\v w')}\to e^{-i\v k\cdot(\v w-\v y')}\,e^{-i\v k\cdot(\v z-\v y')\,\xi}\,,\qquad\xi=\frac{p^+}{k^+}\,,
\end{align}
and $\Kv^m(w'-z)$ by $\frac{1}{1+\xi}\Kv^m(y'-z)$. The second factor of Eq.~\eqref{eq:ex3_phase} is the recoil phase. It is the only place in which the conjugate amplitude knows about the longitudinal momentum of the second gluon through a phase, and it is what cuts off the soft integral below.

\paragraph{Step 4: longitudinal structure.} After Step~3 the three tensor structures come with the factors $\frac{1}{1+\xi}$, $\frac1\xi$ and $1$ under $\int d\xi$:
\begin{align}\label{eq:ex3_three}
\frac{1}{1+\xi}\,\delta_{k'm}\delta_{ij}-\frac1\xi\,\delta_{jm}\delta_{ik'}-\delta_{im}\delta_{jk'}\,.
\end{align}
They behave differently at the two ends of the integration range $\Lambda/k^+<\xi<\vee/k^+-1$.

\emph{The term $1/\xi$ (soft end point).} It is split as $\frac1\xi=\frac{1}{\xi(1+\xi)}+\frac{1}{1+\xi}$. The first part converges at large $\xi$ and is written with a plus prescription, Eq.~\eqref{eq:plus_xi}; the subtracted term is the value of the whole integrand at $\xi=0$, where the recoil phase is one and $\v y'=\v w'$, times
\begin{align}
\int_{\Lambda/k^+}^{\vee/k^+-1}\frac{d\xi}{\xi(1+\xi)}=\log\frac{k^+}{\Lambda}+\log\Big(1-\frac{k^+}{\vee}\Big)+O\Big(\frac{\Lambda}{k^+}\Big)\,.
\end{align}
This is the entry of the row in the JIMWLK list. The second part, $\frac{1}{1+\xi}$, is combined with the first term of Eq.~\eqref{eq:ex3_three}.

\emph{The terms without decay at large $\xi$ (energetic end point).} The third term of Eq.~\eqref{eq:ex3_three} does not decrease with $\xi$ at all. At large $\xi$ the recoil phase oscillates rapidly unless $\v y'\to\v z$, so that only this region contributes, and there the rest of the integrand can be evaluated at $\v y'=\v z$. Subtracting and adding the integrand with $\Kv(x'-y')\to\Kv(x'-z)$ and $U(y')\to U(z)$, and using in the added term
\begin{align}\label{eq:ex3_FT}
\int_{y'}e^{-i\v k\cdot(\v w-\v y')}\,e^{-i\v k\cdot(\v z-\v y')\,\xi}\,\Kv^m(y'-z)=\frac{1}{1+\xi}\int_{w'}e^{-i\v k\cdot(\v w-\v w')}\,\Kv^m(w'-z)\,,
\end{align}
which follows from the homogeneity of $\Kv$, the added term is $-\int d\xi/(1+\xi)=-\log(\vee/k^+)+O(\Lambda/k^+)$ times a structure in which the measured gluon in the conjugate amplitude is emitted from the position $\v z$ of the second gluon. This is the net entry of the row in the BFKL list: the second gluon is a fast gluon of the projectile that acts as a source for the measured one. The two terms with $\frac{1}{1+\xi}$ produce logarithms of $\vee$ in the same way; these two cancel each other, so that the term just described is the whole BFKL coefficient of the row. The subtracted combination is the instantaneous term; it has no pointwise limit for $\vee\to\infty$ and is defined by Eq.~\eqref{eq:inst_xi} with the $\v y'$ integral done first (Appendix~\ref{sec:fin_inst}). \emph{Nuance:} with the plain split $[\frac1\xi]_+$ instead of $[\frac{1}{\xi(1+\xi)}]_+$ the remainder would still grow logarithmically with $\vee$; the split used here takes out the complete $\log(\vee/k^+)$ of the row.

\paragraph{Step 5: the finite remainder.} What is left is, for $\Lambda\to0$ and $\vee\to\infty$,
\begin{align}\label{eq:ex3_fin}
&\frac{1}{(2\pi)^3}\,\frac{ig^4}{4\pi^4}\,\frac{1}{k^+}\int_0^\infty\frac{d\xi}{2\pi}\int_{w,y',z,y,x,x'}e^{-i\v k\cdot(\v w-\v y')}e^{-i\v k\cdot(\v z-\v y')\xi}\notag\\
&\qquad\times\Big[\mathfrak K^{(y')}_{\rm III}\,\{\dots U(y')\dots\}+\mathfrak K^{(z)}_{\rm III}\,\{\dots U(z)\dots\}\Big]\rho\rho\rho+{\rm c.c.}
\end{align}
with the kernels~\eqref{eq:K-y_rm_III} and~\eqref{eq:K-z_rm_III}: the first bracket of $\mathfrak K^{(y')}_{\rm III}$ is the combination of the terms with $\frac{1}{1+\xi}$, its second term the plus-prescription term, its third term and $\mathfrak K^{(z)}_{\rm III}$ the instantaneous term and its subtraction; the colour factor $\{\dots U(y')\dots\}$ is Eq.~\eqref{eq:ex3_colour} and $\{\dots U(z)\dots\}$ the same with $U(y')\to U(z)$.

\paragraph{Step 6: ordered charges and the average.} The ordered product $\rho^a(x')\rho^b(x)\rho^c(y)$ is written, by Eq.~\eqref{eq:symord}, as its symmetrised product plus contact terms. The first gives Eq.~\eqref{eq:H-III.V-3}, the second Eq.~\eqref{eq:H-III.V-2}; the average over the projectile, with Eqs.~\eqref{eq:cubicavg} and~\eqref{eq:traceid}, turns them into Eq.~\eqref{eq:S-III.V-3}, proportional to $\sigma(x',x,y)$ and to a C-odd trace of Wilson lines, and Eq.~\eqref{eq:S-III.V-2}, proportional to $n$.

\paragraph{Step 7: what the row needs from another row.} The plus-prescription term of Eq.~\eqref{eq:ex3_fin} has, at large $|\v z|$, a soft remainder with a tail $\propto\log|\v z|/\v z^2$ that is not integrable (Section~\ref{sec:nuance_soft}). It is cancelled by the term of Row~I.III with the kernel $\mathfrak K^{(w)}_{\rm I.III}$, in which the same soft gluon is absorbed by the measured gluon in the amplitude instead of being emitted in the conjugate amplitude (pair~A of Table~\ref{tab:pairs}). The finite term of Row~III.V is therefore defined only together with that term.

\paragraph{What the example shows.} A single interference row contributes to both evolutions and to the finite terms; the split between them is made by the behaviour in $\xi$ of three different tensor structures; the recoil phase plays the role of a dynamical cut-off; and the finite term is neither a function nor separately integrable, but a distribution in $\xi$ that has to be combined with a virtual term.

\begin{summarybox}
\begin{itemize}
\item \emph{Row I.II, one-$\rho$ part:} LO $\times\frac{\alpha_sN_c}{4\pi}\{(\frac{11}{3}-2L)[\frac4\epsilon+\mathcal B+\mathcal B']+2(\frac{67}{9}-\frac{\pi^2}{3})\}$, Eq.~\eqref{eq:ex1_res}. The numbers $\frac{11}{6}$, $\frac{67}{18}$ and $\frac{\pi^2}{3}$ are the three integrals~\eqref{eq:ex1_ints}.
\item \emph{Row II.IV:} the polarisation sum of two light-cone splitting vertices gives the real splitting function; the collinear pole comes from large separation of the daughters; the end points $\xi\to1$ and $\xi\to0$ give a JIMWLK and a BFKL term; the plus prescription on $\frac{1-\xi}{\xi}$ leaves the remnant, Eqs.~\eqref{eq:ex2_split}, \eqref{eq:ex2_rem}.
\item \emph{Row III.V:} the recoil phase~\eqref{eq:ex3_phase}; the split $\frac1\xi=\frac{1}{\xi(1+\xi)}+\frac{1}{1+\xi}$; the instantaneous term and its subtraction; symmetrised and contact parts; pairing with Row~I.III.
\item The intermediate algebra of this section (longitudinal integrals, polarisation sums, the $d_\perp$-dimensional contraction, prefactors) was re-checked symbolically and numerically; the script is part of the material available from the authors.
\end{itemize}
\end{summarybox}
